\chapter{Type Inference, Algebraic Data Types, and Pattern Compilation}
\label{ch:functional-types-patterns}

An unannotated identity function can be useful at both integer and Boolean arguments.
An unannotated function parameter cannot necessarily be used polymorphically in the same body.
The compiler must distinguish those cases before erasing types or compiling pattern matches.
Inference is not permission to make every use independently acquire whatever type it needs.

\section{Infer monotypes and generalize selected bindings}
\label{sec:ch90-schemes}
The lab monotypes are Int, Bool, type variables, arrows, and List(T).
A \term{type scheme} quantifies selected type variables, written $\forall\alpha_1\ldots\alpha_n.\tau$.
Using a scheme instantiates its quantified variables with fresh unknowns.
Using a monomorphic binding shares the same unknown with every use until constraints resolve it.

The let-bound identity receives scheme $\forall\alpha.\alpha\to\alpha$.
Its application to true instantiates $\alpha$ as Bool; a separate application to seventeen instantiates a fresh variable as Int.
By contrast, a lambda-bound f is monomorphic within its body in this Hindley--Milner fragment.
Requiring both f(true) and f(1) forces an incompatible domain and is rejected.
Table~\ref{tab:ch90-inference} records the cases executed by the tests.

\begin{table}[htbp]
\caption{Fresh instantiation belongs to quantified schemes, not every variable use.}
\label{tab:ch90-inference}
\centering
\begin{tabular}{ll}
\toprule
Program shape & Inference outcome \\
\midrule
let id = lambda x.x; use at Bool and Int & accepted with fresh instances \\
lambda f; require f at Bool and Int domains & incompatible monomorphic constraints \\
lambda x; apply x to x & occurs-check failure \\
cons integer onto Boolean list & element-type mismatch \\
if Boolean then integer else Boolean & branch-type mismatch \\
\bottomrule
\end{tabular}
\end{table}

Generalization quantifies free variables of the inferred binding type that are not free in the current environment.
The environment must be viewed through the current substitution.
Generalizing an unknown already tied to an outer binding would incorrectly allow later uses to escape that shared constraint.
The companion explicitly subtracts environment free variables before forming the scheme.

\section{Generate and solve equality constraints}
\label{sec:ch90-unification}
Application infers function type F and argument type A, creates fresh result variable B, and unifies F with $A\to B$.
A conditional requires a Bool test and equal branch types.
Integer primitives require Int operands, and equality in this lab compares integers only.
An empty list receives List($\alpha$) with a fresh element unknown; cons unifies its tail with List of the head's type.

Algorithm~\ref{alg:ch90-infer} gives let generalization in the incremental-substitution style used by the lab.
It mutates a solver substitution rather than returning an explicit composed substitution at every recursive call.
All subsequent free-variable and type reads resolve through that substitution.
This is the same semantic discipline as Algorithm W-style inference, organized for a small implementation rather than claimed to be a literal presentation of W.

\begin{algorithm}[htbp]
\caption{Infer a nonrecursive let with environment-aware generalization.}
\label{alg:ch90-infer}
\begin{pseudocode}
\Require pure expression let x = value in body; scheme environment; current solver
\Ensure body type or justified type error
\State infer value type $\tau$ while updating equality constraints
\State resolve $\tau$ and every relevant environment type under current substitution
\State $Q\gets\operatorname{FTV}(\tau)\setminus\operatorname{FTV}(environment)$
\State extend environment with $x:\forall Q.\tau$
\State infer body under the extended environment
\end{pseudocode}
\end{algorithm}

The solver uses an occurs check: solving $\alpha=\alpha\to\beta$ would require an infinite finite-type tree and fails.
It distinguishes Bool from Int even though Python's Boolean values are integer subclasses; the source language's type system does not inherit that host-language quirk.
The implementation checks Boolean syntax and types before integer syntax and types.
Its tests include applying a number as a function and mixing list element types.

There is no general linear source-size bound for this elementary inference engine.
Repeated resolution, structural types, free-variable scans, and explicit type growth contribute cost, and principal types can be large.
A production implementation uses sharing, union-find-like structures, levels for generalization, and carefully bounded diagnostic rendering.
A union-find component alone does not eliminate occurs checking or the cost of a very large inferred type.

\section{Handle recursion without polymorphic recursion guesses}
\label{sec:ch90-recursive}
For letrec f=lambda... in body, introduce a fresh monotype for f before inferring its right side.
Infer the lambda while recursive references use that same monotype, then unify the placeholder with the inferred type.
Afterward, generalize relative to the outer environment for uses in the body.
This supports ordinary monomorphic recursion with a polymorphic exported binding when justified.

It does not infer arbitrary polymorphic recursion where recursive calls instantiate the function at structurally different types within its own definition.
Adding such a feature typically needs annotations or a different inference contract.
Likewise, unrestricted let generalization in a language with mutable references requires an effect-sensitive restriction such as a value restriction.
Our core is pure and has no reference allocation, so that additional language feature is not silently assumed.

Figure~\ref{fig:ch90-flow} shows how constraints precede representation choices.
The compiler consumes solved types for primitive legality and constructor shape but the reference runtime retains ordinary tagged host values.
The implementation does not claim unboxed polymorphic specialization merely because its checker inferred a type.

\begin{figure}[htbp]
\centering
\begin{tikzpicture}[node distance=12mm]
\node[basic block] (syntax) {validated core expression};
\node[basic block,below=of syntax] (constraints) {fresh variables and equality constraints};
\node[basic block,below=of constraints] (schemes) {solved types and let schemes};
\node[basic block,below=of schemes] (lower) {typed primitives and constructor lowering};
\draw[cfg edge] (syntax)--(constraints);\draw[cfg edge] (constraints)--(schemes);\draw[cfg edge] (schemes)--(lower);
\end{tikzpicture}
\caption{Inference justifies later operations but does not by itself implement a specialized representation.}
\label{fig:ch90-flow}
\end{figure}

\section{Compile algebraic cases through constructor tests}
\label{sec:ch90-patterns}
An algebraic data type is described by constructors and their field types.
The lab implements the list type with Nil and Cons, while the general design extends to arbitrary declared sums and products.
A case first identifies the constructor, then exposes only that constructor's fields to its branch.
The empty branch and cons branch must have the same result type.

Algorithm~\ref{alg:ch90-case} compiles the lab's exhaustive two-constructor case.
Its dispatch is constant-time once the scrutinee's outer constructor is available; evaluating that scrutinee can take arbitrary computation.
In the lazy machine, exposing Cons does not force its head or tail.
In the strict machine those fields were evaluated during construction.

\begin{algorithm}[htbp]
\caption{Lower a typed list case without evaluating the unselected branch.}
\label{alg:ch90-case}
\begin{pseudocode}
\Require list scrutinee, empty branch, cons field binders and branch; equal result types
\Ensure selected branch with correctly bound fields
\State evaluate scrutinee to its required outer constructor
\If{constructor is Nil} \State execute empty branch
\Else
\State bind head and tail fields in a fresh lexical environment
\State execute cons branch under that environment
\EndIf
\end{pseudocode}
\end{algorithm}

General nested patterns form a matrix of constructors, variables, literals, and wildcards.
Decision-tree compilation chooses a column to inspect, specializes constructor rows, and carries default rows into compatible branches.
Source rule order matters when patterns overlap, and guards can fail after a structural match.
A DAG can share tests across branches, but sharing must not reorder effects or demand lazy fields earlier than the source semantics allows.

Figure~\ref{fig:ch90-case} gives the simple list decision tree.
Coverage checking for general nested and recursive types is more involved than counting distinct first-level constructors.
Redundancy checking asks whether an earlier pattern/guard combination already covers a later one.
The lab validates only its fixed two-arm syntax and does not implement a general pattern-matrix compiler or exhaustiveness solver.

\begin{figure}[htbp]
\centering
\begin{tikzpicture}[node distance=14mm and 15mm]
\node[basic block] (tag) {inspect outer list constructor};
\node[basic block,below left=of tag] (nil) {Nil: empty branch};
\node[basic block,below right=of tag] (cons) {Cons: bind h,t; body};
\draw[cfg edge] (tag)--(nil);\draw[cfg edge] (tag)--(cons);
\end{tikzpicture}
\caption{Constructor dispatch exposes fields only on the branch where their representation exists.}
\label{fig:ch90-case}
\end{figure}

\section{Use types without discarding diagnostics}
\label{sec:ch90-diagnostics}
A failed unification is most useful when it names the two conflicting constraints and their source spans.
The minimal lab reports mismatch categories but does not retain rich source spans through inference.
A production elaborator should preserve provenance while normalizing types, especially when one occurrence creates a constraint and a distant occurrence reveals its inconsistency.
Pretty-printing inferred type variables deterministically also makes regression tests and user reports reproducible.

Table~\ref{tab:ch90-pattern-obligations} lists contracts that richer language additions would require.
Type erasure preserves behavior only when runtime operations still implement the checked representation and required dynamic tests.
An optimizer replacing a typed node must update or recompute its attached type facts, not retain a stale annotation from the old expression.

\begin{table}[htbp]
\caption{Richer functional-language features add static and runtime obligations together.}
\label{tab:ch90-pattern-obligations}
\centering
\begin{tabular}{ll}
\toprule
Feature & Additional machinery \\
\midrule
general algebraic types & declarations, constructor schemes, layout metadata \\
nested patterns & ordered decision compilation and coverage checks \\
mutable references & restricted generalization and effect-aware optimization \\
polymorphic recursion & explicit annotation/inference policy \\
unboxed specialization & calling and representation consistency at every use \\
\bottomrule
\end{tabular}
\end{table}

\section{Work a recursive list function's constraints}
\label{sec:ch90-sum-constraints}
The sum example illustrates how several local constraints determine a recursive type together.
Give the recursive binding sum a fresh placeholder $\alpha$ and the lambda parameter xs a fresh type $\beta$.
The case scrutinee requires $\beta=\operatorname{List}(\gamma)$ for fresh element type $\gamma$.
The empty branch returns integer zero, so the case result must be Int.
In the cons branch, head has type $\gamma$ and tail has type List($\gamma$).

The addition requires head to be Int, solving $\gamma=\operatorname{Int}$.
The recursive call sum(tail) constrains $\alpha$ to an arrow from List(Int) to a fresh result type, and addition requires that result to be Int.
The lambda therefore has type List(Int)$\to$Int.
Unifying it with the recursive placeholder is consistent and closes the recursive obligation.
No polymorphic recursive instantiation was needed inside the body.

This example also explains why constraints must flow both forward and backward through syntax.
Nil alone leaves its element type unknown, but a later cons head or case branch can constrain it.
Inferring each subexpression once and freezing its first printed type would miss those equations.
The solver keeps references to shared unknowns and resolves them under the current substitution whenever a type is consumed or generalized.
After successful inference, the case compiler can trust that both arms return Int and that the cons fields have the promised element and list types.

\section{Exercises}
\label{sec:ch90-exercises}
\begin{enumerate}
\item\label{ex:ch90-id} Derive the scheme of let-bound identity and its two instances for true and seventeen.
\item\label{ex:ch90-self} Generate the type equation for lambda x.app(x,x) and explain the occurs-check failure.
\item\label{ex:ch90-list} Infer the type of cons(4,cons(7,nil)) and locate the conflict if seven is replaced by true.
\item\label{ex:ch90-case} Type a case returning eight for Nil and the head for Cons.
What type must its list elements have?
\end{enumerate}

\section{Further reading}
\label{sec:ch90-reading}
Milner develops the polymorphic typing foundation~\cite{milner-1978-polymorphism}.
Appel develops ML implementation, inference, and pattern compilation in a larger compiler setting~\cite{appel-1998-modern}.
The finite-term unification chapter provides the occurs-check foundation, while the earlier type-theory chapter distinguishes explicit checking from inference.
